\begin{proof}
    Clearly $\Sigma_{\textbf{Prim}}$, $\Sigma_{\textbf{Ctrl}}$, and $\Sigma_{\textbf{Pow}}$ are pairwise disjoint.
    This is because $( \odot )$ and $( \exparr )$ are distinct term constructors and $\Sigma_{\textbf{Prim}}$ consists only of base terms.
    This means that the set $\Sigma_{\textbf{Prim}} \sqcup \Sigma_{\textbf{Ctrl}} \sqcup \Sigma_{\textbf{Pow}}$ is well-defined.
    Moreover, $\Sigma_{\textbf{Prim}}$, $\Sigma_{\textbf{Ctrl}}$, and $\Sigma_{\textbf{Pow}}$ are all subsets of $\Sigma_{\CatHQC}$ by construction.
    Then it suffices to show that $\Sigma_{\CatHQC}$ is a subset of $\Sigma_{\textbf{Prim}} \sqcup \Sigma_{\textbf{Ctrl}} \sqcup \Sigma_{\textbf{Pow}}$.
    Let $g \in \Sigma_{\CatHQC}$.
    Then there exists some least $n \in \mathbb{N}$ such that $g \in \Sigma_{\CatHQC}^n$.
    If $n = 0$, then it follows that $g \in \Sigma_{\CatHQC}^0 = \Sigma_{\textbf{Prim}}$.
    Otherwise, $n > 0$ and $g \in \Sigma_{\CatHQC}^n \setminus \Sigma_{\CatHQC}^{n-1}$.
    By definition, the following equation holds.
    \begin{align*}
        \Sigma_{\CatHQC}^n \setminus_{\CatHQC}^{n-1}
        &=
        \left(
            \Sigma_{\CatHQC}^{n-1}
            \cup
            \Tower( \CatHQC[n-1] )
            \cup
            \Power( \CatHQC[n-1] )
        \right)
        \setminus
        \Sigma_{\CatHQC}^{n-1}
        \\
        &\subseteq
        \Tower( \CatHQC[n-1] )
        \cup
        \Power( \CatHQC[n-1] )
    \end{align*}
    Then either $g \in \Tower( \CatHQC[n-1] )$ or $g \in \Power( \CatHQC[n-1] )$.
    There are three cases to consider.
    \begin{enumerate}
    \item Assume that $g \in \Tower( \CatHQC[n-1] )$ and $|g| = 2$
          Then $g_1 \in \Mor( \CatHQC[n-1] ) \subseteq \CatHQC$ and $g_2 \in \Mor( \CatHQC[n-1] ) \subseteq \CatHQC$.
          Moreover, $g = g_1 \odot g_2$.
          In conclusion, $g \in \Sigma_{\mathbf{Ctrl}}$.
    \item Assume that $g \in \Tower( \CatHQC[n-1] )$ and $|g| > 2$
          Then $g_1 \in \Mor( \CatHQC[n-1] ) \subseteq \CatHQC$.
          Define $V = g_2 \odot g_3 \odot \cdots \odot g_k$ where $k = |g|$.
          Then the following sequence of inclusions holds.
          \begin{equation*}
              V \in \Tower( \CatHQC[n-1] )
              \subseteq
              \Sigma_{\CatHQC}^n
              \subseteq
              \CatHQC[n]
              \subseteq
              \CatHQC
          \end{equation*}
          Moreover, $g = g_1 \odot V$.
          In conclusion, $g \in \Sigma_{\mathbf{Ctrl}}$.
    \item Assume that $g \in \Power( \CatHQC[n-1] )$.
          Then there exists a $U \in \CatHQC[n-1]$ and $\alpha \in \mathbb{R}$ such that $g = U \exparr \alpha$.
          Since $\CatHQC[n-1] \subseteq \CatHQC$, then $g \in \Sigma_{\mathbf{Pow}}$.
    \end{enumerate}
    In either case, $g \in \Sigma_{\textbf{Ctrl}} \sqcup \Sigma_{\textbf{Pow}}$.
    Then $g \in \Sigma_{\textbf{Prim}} \sqcup \Sigma_{\textbf{Ctrl}} \sqcup \Sigma_{\textbf{Pow}}$ regardless of the value of $n$.
    Since $g$ was arbitrary, then $\Sigma_{\CatHQC} \subseteq \Sigma_{\textbf{Prim}} \sqcup \Sigma_{\textbf{Ctrl}} \sqcup \Sigma_{\textbf{Pow}}$.
    This completes the proof.
\end{proof}